\clearpage
\section{Renaming the literal atomic translations}
\label{sec:copying-atomic-renaming}
\termdef{copying-atomic-renaming}
This section closes the selected-resource and all three atomic cases of
bijective environment renaming for the original Figure~16 translation.
The construction works for arbitrary finite variable and atom environments,
at every adequate allowance for the individual clauses. The proof retains
all weakening and contraction edges and the atom clause's instance edge.
It neither simplifies a distributor to a wire nor replaces the proposition's
terms. Composite-term renaming and general substitution remain separate tasks.

\subsection{Coordinates for the selected-resource graph}
\termdef{copying-selected-resource-coordinates}
For a resource $i<n$, the raw assembly is precisely
$(D_{i,0}\otimes D_{1,1})\otimes D_{n-i-1,0}$, the expression inside the
existing interface cast. We introduce coordinates for its carriers, not a
new translation. There is one weakening for every output resource and one
additional contraction belonging to the selected resource.

\lean{Copying/SelectedResourceShape}{9}{11}{SelectedResourceShape.raw}

The edge-coordinate map sends a weakening in the left forest to its original
resource number, the selected weakening to $i$, and a weakening at offset $j$
in the right forest to $i+1+j$. The selected contraction is the extra summand
of $\operatorname{Fin}(n)+1$. The inverse compares a resource number with $i$:
less than $i$ selects the left forest; equality selects the middle weakening;
greater than $i$ selects the right forest after subtracting $i+1$. The extra
summand selects the contraction. Checking these three cases proves both
inverse identities, including the empty left or right forest cases, and
packages them as an edge bijection.

\lean{Copying/SelectedResourceShape}{13}{16}{SelectedResourceShape.edgeIndex}
\lean{Copying/SelectedResourceShape}{18}{22}{SelectedResourceShape.edgeAt}
\lean{Copying/SelectedResourceShape}{24}{34}{SelectedResourceShape.edgeIndex_at}
\lean{Copying/SelectedResourceShape}{36}{62}{SelectedResourceShape.edgeAt_index}
\lean{Copying/SelectedResourceShape}{64}{65}{SelectedResourceShape.edges}

The internal carrier is a singleton: only the middle one-use tree has an
internal vertex. The two weakening forests have empty internal carriers.
Case analysis on the tensor and forest tags proves the corresponding
bijection with the unit type. The input length is one; the output length is
$i+1+(n-i-1)=n$, using $i<n$. These identities give a bijection from all raw
vertices to $\operatorname{Port}(1,1,n)$: preserve each boundary index and
use the singleton coordinate internally. The inverse laws follow from the
internal bijection and equality of finite indices with the same value.

\lean{Copying/SelectedResourceShape}{68}{84}{SelectedResourceShape.internal}
\lean{Copying/SelectedResourceShape}{86}{88}{SelectedResourceShape.input_length}
\lean{Copying/SelectedResourceShape}{90}{94}{SelectedResourceShape.output_length}
\lean{Copying/SelectedResourceShape}{96}{113}{SelectedResourceShape.ports}

In these coordinates every weakening has empty source. The contraction's
source is the ordered pair consisting first of the internal vertex and then
of the sole input. Weakening $i$ targets the internal vertex; every other
weakening $j$ targets output $j$; the contraction targets output $i$.
The label-coordinate function assigns weakening to the resource summand
and contraction to the extra summand.

\lean{Copying/SelectedResourceShape}{116}{118}{SelectedResourceShape.source}
\lean{Copying/SelectedResourceShape}{120}{122}{SelectedResourceShape.target}
\lean{Copying/SelectedResourceShape}{124}{126}{SelectedResourceShape.label}

The source, target and label correspondence lemmas prove these descriptions
for the actual raw assembly, by splitting the resource comparison and the
two middle edges. The target proof checks the right-forest offset arithmetic
and the equality test for the selected resource. All vertices have the
selected wire sort: the sole internal case is direct, and boundary cases
use membership in the three concatenated replicated interface lists.
Thus coordinates describe the entire ordered incidence data, not merely
edge counts or an informal shape.

\lean{Copying/SelectedResourceShape}{130}{143}{SelectedResourceShape.source_at}
\lean{Copying/SelectedResourceShape}{145}{162}{SelectedResourceShape.target_at}
\lean{Copying/SelectedResourceShape}{166}{173}{SelectedResourceShape.label_at}
\lean{Copying/SelectedResourceShape}{175}{192}{SelectedResourceShape.vertex_label}

\subsection{Renaming the original selected resource}
For a resource bijection $\pi$, permute output coordinates by $\pi$ and fix
the internal vertex and input. On edge coordinates apply $\pi$ to weakening
resources and fix the extra contraction. The source lists commute with these
maps directly. Target lists commute because
$\pi(j)=\pi(i)$ if and only if $j=i$. This uses injectivity and is the point
where the argument is specifically bijective renaming.

\lean{Copying/SelectedResourceRenaming}{8}{12}{SelectedResourceRenaming.outputPermutation}
\lean{Copying/SelectedResourceRenaming}{14}{16}{SelectedResourceRenaming.source_permutation}
\lean{Copying/SelectedResourceRenaming}{18}{25}{SelectedResourceRenaming.target_permutation}

Applying the inverse port-coordinate bijection to the source and target
correspondence lemmas reconstructs the actual incidence lists. Compose the
source coordinate bijection, the permutation, and the inverse target
coordinate bijection, on vertices and edges alike. Wire labels agree because
all these vertices have the same wire sort; input and output roles are
preserved. The coordinate label lemma gives recursive equality of edge
labels. Reconstruct the source and target lists and use their permutation
identities to finish the boundary-isomorphism proof.
The raw output-action lemma states explicitly that output $j$ is sent to
output $\pi(j)$.

\lean{Copying/SelectedResourceRenaming}{27}{30}{SelectedResourceRenaming.source_eq}
\lean{Copying/SelectedResourceRenaming}{32}{35}{SelectedResourceRenaming.target_eq}
\lean{Copying/SelectedResourceRenaming}{38}{68}{SelectedResourceRenaming.rawIso}
\lean{Copying/SelectedResourceRenaming}{69}{73}{SelectedResourceRenaming.rawIso_output}

Interface casts transport boundary isomorphisms by equality elimination.
A companion lemma transports a specified output action along the induced
finite-index casts. Applying these to the raw isomorphism gives the theorem
for the existing public selected-resource constructor, and proves that its
output action is exactly $j\mapsto\pi(j)$. This checks the boundary action
of the actual Figure~16 clause after its original casts.

\lean{Copying/SelectedResourceRenaming}{81}{84}{SelectedResourceRenaming.castInterfaces}
\lean{Copying/SelectedResourceRenaming}{85}{90}{SelectedResourceRenaming.castInterfaces_output}
\lean{Copying/SelectedResourceRenaming}{96}{100}{SelectedResourceRenaming.selectedResource_rename}
\lean{Copying/SelectedResourceRenaming}{102}{108}{SelectedResourceRenaming.selectedResource_rename_output}

\subsection{Transport through ordered parallel assembly}
For a family of component boundary isomorphisms, map tagged vertices and edges
componentwise, retaining their component tags. Each inverse identity reduces
to the component inverse identity. Compose the inverse original parallel
vertex assembly with this tagged vertex map and the target assembly to get
the whole-graph port bijection. The embedding lemma says that this bijection
commutes with embedding each component vertex.

\lean{Copying/BoundaryParallel}{11}{15}{BoundaryParallel.componentPorts}
\lean{Copying/BoundaryParallel}{17}{24}{BoundaryParallel.componentEdges}
\lean{Copying/BoundaryParallel}{26}{28}{BoundaryParallel.ports}
\lean{Copying/BoundaryParallel}{30}{35}{BoundaryParallel.embed}

The resulting assembly isomorphism preserves wire labels using the two
embedding-label equations and the component isomorphism. The input and output
role proofs reduce through the embedding equations to the corresponding
component role proofs. Edge labels are related componentwise. For sources
and targets, map the component incidence equality through the embedding and
use the embedding lemma entry by entry. Specialising to the original Boolean
family and concatenated boundary gives ordered tensor congruence.

\lean{Copying/BoundaryParallel}{38}{79}{BoundaryParallel.iso}
\lean{Copying/BoundaryParallel}{85}{90}{BoundaryParallel.tensor}

The variable clause is now immediate from selected-resource renaming on the
term side and the zero-use atom-distributor automorphism on the atom side,
followed by the original interface casts. For a constant, tensor the zero-use
term-distributor automorphism, the identity of the constant operation, and
the zero-use atom-distributor automorphism, then use the original casts.
Unused resources remain represented by their weakening edges.

\lean{Copying/BoundaryParallel}{96}{101}{BoundaryParallel.variableTerm_rename}
\lean{Copying/BoundaryParallel}{104}{110}{BoundaryParallel.constant_rename}

\subsection{Retaining the atom's instance edge}
The atom clause additionally has the literal serial prefix of an instance
edge. Let $H$ be its remaining graph, with one atom input. The prefixed graph
has one new external term input; every old vertex survives, with the old
atom input becoming internal. A bijection between its vertices and
$1+\operatorname{Vertex}(H)$ records this fact, checking both depth-zero and
positive-depth representations. Lift a boundary isomorphism of $H$ by fixing
the new input and applying its old port bijection.

\lean{Copying/BoundaryInstance}{8}{8}{BoundaryInstance.graph}
\lean{Copying/BoundaryInstance}{10}{35}{BoundaryInstance.ports}
\lean{Copying/BoundaryInstance}{37}{39}{BoundaryInstance.liftPorts}

A boundary isomorphism between graphs with one input fixes that input:
input-role preservation implies its image is an input, and the index set
has only one element. Embedding an old vertex uses the inverse coordinate
map on the right summand. The lifted map commutes with this embedding by
cancellation. Embedded vertices retain their wire labels; none is a new
external input; and being an output is preserved and reflected. These three
facts follow by cases on the old vertex. The actual expansion insertion
agrees with this embedding, including the shared internal atom vertex and
the original ordered output positions.

\lean{Copying/BoundaryInstance}{41}{47}{BoundaryInstance.input_fixed}
\lean{Copying/BoundaryInstance}{50}{51}{BoundaryInstance.embed}
\lean{Copying/BoundaryInstance}{53}{56}{BoundaryInstance.lift_embed}
\lean{Copying/BoundaryInstance}{60}{62}{BoundaryInstance.vertex_embed}
\lean{Copying/BoundaryInstance}{64}{65}{BoundaryInstance.input_embed}
\lean{Copying/BoundaryInstance}{67}{68}{BoundaryInstance.output_embed}
\lean{Copying/BoundaryInstance}{70}{73}{BoundaryInstance.insert}
\lean{Copying/BoundaryInstance}{75}{86}{BoundaryInstance.insert_eq_embed}

Map the new instance edge to itself and the old edges by the given edge
bijection. Its inverse identities hold separately in the two tags. For the
prefix isomorphism, the vertex and role laws use the preceding coordinate
facts. The new edge's source is fixed; its target is fixed by uniqueness of
the old atom input. For every old edge, source and target preservation follow
from the old incidence equations, insertion/embedding correspondence, and
commutation of the lifted port map with embedding. Tensor this prefix result
with the ambient term-weakening forest, and apply the original casts, to
obtain renaming for the public atom translation.

\lean{Copying/BoundaryInstance}{88}{95}{BoundaryInstance.edges}
\lean{Copying/BoundaryInstance}{97}{159}{BoundaryInstance.iso}
\lean{Copying/BoundaryInstance}{165}{171}{BoundaryInstance.atom_rename}

\subsection{The atomic-case endpoint and remaining work}
The predicate \texttt{Atomic} selects precisely variables, atoms and constants
from the original scoped syntax. The atomic renaming theorem proves existence
of a boundary isomorphism between the original translation and the translation
of the syntactically renamed term. Its sole depth cast uses the already proved
depth-renaming identity. Cases on the atomic predicate reduce to the three
public clause results above. There are no operational or template hypotheses.

\lean{Copying/AtomicRenaming}{8}{14}{AtomicRenaming.Atomic}
\lean{Copying/AtomicRenaming}{15}{22}{AtomicRenaming.atomic_rename}

This completes the selected-resource and atomic cases of bijective translation
renaming. It does not prove the inductive cases for application, unary and
binary operations, or lambda boxes; nor does it prove non-bijective renaming,
weakening or graph-level substitution. Boundary-permuting isomorphism is
still distinct from the boundary-fixing equivalence used for observations.
The original terms and translation of Proposition~12.2 are unchanged.

\paragraph{Completion estimate.}
As of 25 October 2026, the estimated completion of the requested
Proposition~12.2 development is about \textbf{40\% by mathematical effort},
with a plausible range of \textbf{30--50\%}. This is a judgement about remaining
proof difficulty, not a measured percentage of declarations or checklist items.
The current milestone closes actual translation cases; it does not erase the
large remaining concrete copying-machine and eleven-template obligations.
